% !TeX root = ../main.tex

\newlecture[Present Value and Internal Rate of Return][Sep 17, 2026]
Present value: comparing cash flows that occur at different times, while interest rates tells how to move cash flows in time. Use interest rates to move all cash flows to the same time, so they become easy to compare.

\b{Ideal Bank}: applies the same rate of interest to deposits and loans, and has no transaction costs (real banks have higher rates for loans...).
Used for two ways:
\begin{itemize}
    \item \rb{Compounding}: when cash flow is moved \b{forward}, \dbb{multiplied} by $(1+r)$ each period.
    \item \rb{Discounting}: when cash flow is moved \b{backward}, \dbb{divided} by $(1+r)$ each period.
\end{itemize}

\header{The Comparison Principle}
The ideal bank brings cash flow stream to the same time, then is compared.
\begin{itemize}
    \item When all cash flows are moved to \textit{present time}, it is called \rb{present value}.
    \item When all cash flows are moved to the \textit{same future timeline}, it is called \rb{future value}.
\end{itemize}

\begin{example}
    Cash flow stream $(-2,1,1,1)$ at $r = 10\%$ (yearly compounding, but any compounding convention works):
    \begin{align*}
        FV &= -2(1.1)^3 + 1(1.1)^2 + 1(1.1) + 1 = 0.648\\
        PV &= -2 + \frac{1}{1.1} + \frac{1}{(1.1)^2} + \frac{1}{(1.1)^3} = 0.487
    \end{align*}
    Check: $PV = FV/(1.1)^3 = 0.648/(1.1)^3 = 0.487$
\end{example}

\header{General Formula for Present Value}
Let $x = (x_0, x_1, x_2, \dots, x_n)$
\begin{align*}
    PV(x) &= x_0 + \frac{x_1}{(1+r)} + \frac{x_2}{(1+r)^2} + \cdots + \frac{x_n}{(1+r)^n} \\
    &= \boxed{\sum_{k=0}^{n}\frac{x_k}{(1+r)^k}}
\end{align*}
Note that it's usually assumed that $x_0$ is negative, for the initial investment.


\header{Main Theorem on Present Value}
Two cash flow streams are equivalent for an ideal bank iff \underline{they have the same present value}.

\header{Internal Rate of Return}
Captures the idea on the "return" on an investment.
\\The return on your investment is the return which would make the present value equal to 0.
\begin{example}
    Invest \$100 today, receive \$110 in one year. What is your return?
    \\Intuitively: 10\%, why?
    \[x(-100, 110)\]
    \[PV(x) = -100 + \frac{110}{1+0.1} = 0\]
\end{example}

\begin{definition}
    Let $x = (x_0, x_1, \dots, x_n)$ be a cash flow stream. \\Then, the \rb{internal rate of return} of this stream is the number $r$ satisfying:
    \[\boxed{\begin{gathered}
        PV(x) = 0 \\
        x_0 + \frac{x_1}{1+r} + \frac{x_2}{(1+r)^2} + \cdots + \frac{x_n}{(1+r)^n} = 0
    \end{gathered}}\]
    To solve for this $r$, its more convenient to set 
    \[c = \frac{1}{(1+r)}\]
    and solve
    \[x_0 + x_1c+ x_2 c^2 + \cdots + x_n c^n = 0\]
    with the fundamental theory of algebra.
\end{definition}

\begin{example}
    What is the internal rate of return of $(-2,1,1,1)$?
    \\Solve: $0 = -2 + c + c^2 + c^3$
    \[c = 0.81 \quad \Rightarrow \quad c = \frac{1}{1+r} \quad \Rightarrow \quad r = 0.23\]
\end{example}
The polynomial can have many roots, which can be a problem. Fortunately, most standard investments have a well-defined IRR (see the theorem below).


\heading{Main Theorem on IRR}
Suppose the cash flow stream $x = (x_0, x_1, \dots, x_n)$ has $x_0 < 0$ and $x_k \geq 0$ where at least one term is strictly positive. Then, there is a unique positive root to the equation
\[x_0+x_1c + x_2 c^2 + \dots + x_n c^n = 0\]
Proof: plot $f(c) = x_0 + x_1c + x_2c^2 + \dots + x_n c^n$

\begin{center}
\begin{tikzpicture}
    \draw[->] (-0.3,0) -- (4.5,0) node[right] {$c$};
    \draw[->] (0,-1.5) -- (0,1.8) node[above] {$f(c)$};
    \draw[thick, blue, domain=0:4, smooth, variable=\x] plot ({\x}, {-1 + 0.15*\x*\x});
    \draw[dashed] (2.58,0) -- (2.58,-1.5);
    \filldraw (2.58,0) circle (1.5pt);
    \node[below] at (2.58,-0.05) {$c^*$};
    \filldraw (0,-1) circle (1.5pt);
    \node[left] at (0,-1) {$x_0$};
\end{tikzpicture}
\end{center}
Since $f(0) = x_0 < 0$ and $f'(c) = x_1 + 2x_2c + \dots + nx_nc^{n-1} \geq 0$ for $c \geq 0$ (as each $x_k \geq 0$, with at least one strictly positive), $f$ is nondecreasing and unbounded above. So $f$ crosses zero exactly once, at some $c^* > 0$.

\heading{Net Present Value (NPV)}
Computes all present values of \textit{all} cash flows associated with an investment. Includes the investment term $x_0$ as well.

\begin{example}
    Given two cash flow vectors
    \begin{enumerate}[label=(\arabic*)]
        \item After 1 year: $(-1,2)$
        \item After 2 years: $(-1,0, 3)$
        \\Using NPV solution (choose largest NPV):
        \[-1 + \frac{2}{1.1} = 0.82 \quad \text{vs} \quad -1 + \frac{0}{1.1} + \frac{3}{1.1^2} = 1.48\]
        NPV says choose option (2).
        \\Using IRR (choose greatest IRR):
        \[-1 + 2c = 0 \Rightarrow c = 0.5 \Rightarrow r = 1.0 \quad \text{vs} \quad -1+3c^2 = 0 \Rightarrow c = 0.577 \Rightarrow r = 0.7\]
        IRR says choose option (1).
    \end{enumerate}
\end{example}

\underline{IRR and NPV give conflicting recommendations.}
\\When scaled, \underline{IRR is scale-invariant, PV is not}. 
\\PV also \underline{distinguishes between single and repeated cash flows}, IRR does not.
\\In general, NPV is preferred.

\begin{example}
    Scaling and repetition (PV at 10\%):
    \begin{center}
    \begin{tabular}{lccc}
        \toprule
        Cash flows & $(-1,2)$ & $10\cdot(-1,2)$ & $(-1,2,-1,2,\dots)$\\
        \midrule
        PV & 0.818 & 8.18 & 4.71\\
        IRR & $r=1.0$ & $r=1.0$ & $r=1.0$\\
        \bottomrule
    \end{tabular}
    \end{center}
\end{example}

NPV is preferred because it distinguishes between more cash flows, but this means the cash flows and assumptions need to be modelled carefully before the NPV analysis is done. E.g. the tree-cutting NPV analysis assumes it is done only \textit{once}; if the process can be repeated many times, we arrive at a different answer. Methods of IRR analysis exist that make it agree with NPV (e.g. incremental IRR analysis).

\heading{Steps to Solving Cash Flow Problems}
\begin{enumerate}
    \item Write down the cash flow diagram for each alternative (include \textit{all} cash flows).
    \item Evaluate the NPV of each cash flow.
    \item Choose the alternative with the largest NPV.
\end{enumerate}

\heading{Cycle Problems}
Relate to activities that can be repeated. With NPV you arrive at a different answer depending on how many times an activity is repeated, so the activities must be evaluated over the \underline{same time horizon}. Two approaches:
\begin{enumerate}
    \item Repeat both activities until they terminate at the same time.
    \item Repeat both to infinity: for an activity with lifetime $T$ and one-cycle present value $PV_1$,
    \[PV_1\sum_{k=0}^{\infty}\frac{1}{(1+r)^{Tk}} = PV_1\cdot\frac{1}{1-\frac{1}{(1+r)^T}}\]
\end{enumerate}

\begin{example}[Refrigerator Purchase]
    Model A: price \$600, yearly cost \$180 (beginning of year), life 3 years.
    \\Model B: price \$800, yearly cost \$120 (beginning of year), life 5 years.
    \\Interest rate is 10\%. One cycle of each (price and first cost at time 0):
    \[PV_A = 780 + 180\sum_{k=1}^{2}\frac{1}{1.1^k} = \$1,092.40 \qquad PV_B = 920 + 120\sum_{k=1}^{4}\frac{1}{1.1^k} = \$1,300.38\]
    \b{Approach 1}: equal horizons, 15 years = 5 cycles of A = 3 cycles of B
    \begin{align*}
        A&: PV_A\left(1+\frac{1}{1.1^3}+\frac{1}{1.1^6}+\frac{1}{1.1^9}+\frac{1}{1.1^{12}}\right) = \$3,341.11\\
        B&: PV_B\left(1+\frac{1}{1.1^5}+\frac{1}{1.1^{10}}\right) = \$2,609.17
    \end{align*}
    \b{Approach 2}: repeat to infinity
    \begin{align*}
        A&: 1092.40\cdot\frac{1}{1-\frac{1}{1.1^3}} = \$4,392.70\\
        B&: 1300.38\cdot\frac{1}{1-\frac{1}{1.1^5}} = \$3,430.37
    \end{align*}
    These are costs, so the smaller is better: choose Model B in both approaches.
\end{example}

\heading{Summing Geometric Series}
\begin{align*}
    S = \sum_{k=0}^{\infty}x^k &= 1 + x + x^2 + \cdots\\
    &= 1 + x(1+x+x^2+\cdots) = 1 + xS
\end{align*}
\[S(1-x) = 1 \quad \Rightarrow \quad S = \frac{1}{1-x} \qquad (|x|<1)\]
